\documentclass[11pt]{article}
\usepackage[a4paper,margin=1.9cm]{geometry}
\usepackage[T1]{fontenc}
\usepackage{textcomp}
\usepackage[spanish]{babel}
\usepackage{amsmath,amssymb}
\usepackage{lmodern}
\renewcommand{\familydefault}{\sfdefault}
\usepackage{enumitem}

\newcommand{\vect}[2]{\begin{pmatrix}#1\\#2\end{pmatrix}}
\usepackage{tikz}
\usetikzlibrary{calc,arrows.meta,patterns,decorations.pathmorphing}
\usepackage[table]{xcolor}
\usepackage{multicol}
\usepackage{parskip}

\definecolor{unalblue}{RGB}{31,73,124}

\newlist{opciones}{enumerate}{1}
\setlist[opciones]{label=(\alph*),itemsep=6pt,topsep=6pt,leftmargin=1.8em,font=\normalfont}
\setlist[enumerate,1]{itemsep=1.4em,topsep=1em}

\newcommand{\examfooter}{%
  \vfill
  \rule{\linewidth}{0.4pt}\\[1pt]
  {\scriptsize\textit{Transcripción digital --- MathUNAL}\hfill\textcolor{unalblue}{\textbf{\thepage}}}
}
\pagestyle{empty}

\newcommand{\nota}[1]{{\small\itshape\color{unalblue!70!black} #1}}

\begin{document}
\noindent
\begin{minipage}[t]{0.6\linewidth}
{\small\bfseries Geometría Vectorial y Analítica --- Parcial 1}\\
{\small Semestre 2017--01}
\end{minipage}%
\hfill
\begin{minipage}[t]{0.35\linewidth}
\raggedleft\small Escuela de Matemáticas. Universidad Nacional de Colombia, Sede Medellín.
\end{minipage}\\[2pt]
\rule{\linewidth}{0.6pt}

\noindent
\begin{minipage}[t]{0.54\linewidth}
\textbf{Primer Examen Parcial de Geometría}\\
18 de marzo del 2017\\[6pt]
\textbf{Puntaje. Sólo para uso oficial}

\small
\begin{tabular}{|c|c|c|c|c|c|}
\hline
1 (/30) & 2 (/30) & 3 (/10) & 4-6 (/30) & Total & Nota \\
\hline
 & & & & & \\
\hline
\end{tabular}
\end{minipage}%
\begin{minipage}[t]{0.44\linewidth}
Nombre: \underline{\hspace{4.8cm}}\\[8pt]
Cédula: \underline{\hspace{1.8cm}}\ \ Salón: \underline{\hspace{1.2cm}}\\[8pt]
Grupo: \underline{\hspace{1.8cm}}\ \ Examen No.: \underline{\hspace{1.6cm}}\\[8pt]
Firma: \underline{\hspace{4.8cm}}
\end{minipage}

\vspace{10pt}
\textbf{Instrucciones:} La duración del examen es de 1 hora y 40 minutos. El examen consta de ocho preguntas en tres hojas impresas por ambos lados, antes de empezar verifique que su examen esté \textbf{completo} y que sus \textbf{datos} sean \textbf{correctos}. No \textbf{desengrape} su exámen ni \textbf{añada} otras grapas o su examen será \textbf{anulado}. En las preguntas con procedimiento \textbf{justifique} sus respuestas en los \textbf{espacios asignados}. No está permitido hacer preguntas, el uso de calculadoras, sacar hojas en blanco ni ningún tipo de apuntes durante el examen. Por favor verifique que su celular esté \textbf{apagado}.

\vspace{6pt}
\nota{En cada uno de los literales de los puntos 1, 2 y 3, debe presentar detalladamente el procedimiento para llegar a la solución. NOTA: Respuestas sin procedimiento no se tendrán en cuenta.}

\begin{enumerate}
\item Considere los puntos $A=\vect{-5}{-4}$, $B=\vect{0}{-3}$ y $C=\vect{-6}{0}$.

\begin{enumerate}[label=(\alph*)]
\item{}[7pt] Verifique que los puntos $A$, $B$ y $C$ no son colineales.
\item{}[6pt] Encuentre una ecuación general de la recta que pasa por los puntos $B$ y $C$.
\item{}[6pt] Halle la proyección ortogonal del punto $A$ sobre la recta que pasa por los puntos $B$ y $C$.
\item{}[6pt] Calcule la distancia del punto $A$ a la recta que pasa por los puntos $B$ y $C$.
\end{enumerate}

\item Considere los puntos $A=\vect{-1}{-1}$ y $B=\vect{1}{3}$.

\begin{enumerate}[label=(\alph*)]
\item{}[8pt] Halle una ecuación vectorial de la circunferencia cuyo diámetro es el segmento $\overline{AB}$.
\item{}[8pt] Halle una ecuación vectorial de la mediatriz del segmento $\overline{AB}$.
\item{}[9pt] Suponga que el segmento $\overline{AB}$ es la diagonal de un cuadrado. Halle las coordenadas de los otros dos puntos del cuadrado.
\end{enumerate}

\item{}[10pt] En la figura se observa un rectángulo $\mathcal{R}$ cuyos vértices son los puntos $A$, $B$, $C$, $D$. El punto $M$ es el punto de intersección de las diagonales del rectángulo $\mathcal{R}$. Demuestre la siguiente igualdad
$$A+B+C+D=4M.$$
\textbf{Indicación.} Recuerde que el punto $M$ divide en dos partes iguales a cada una de las diagonales.

\begin{center}
\begin{tikzpicture}[scale=0.9]
\coordinate (B) at (0,2.2);
\coordinate (A) at (3.3,3.5);
\coordinate (D) at (3.9,2.5);
\coordinate (C) at (0.6,1.2);
\coordinate (M) at ($(B)!0.5!(D)$);
\draw[thick] (A) -- (B) -- (C) -- (D) -- cycle;
\draw[dashed] (A) -- (C);
\draw[dashed] (B) -- (D);
\node[above right] at (A) {$A$};
\node[left] at (B) {$B$};
\node[below left] at (C) {$C$};
\node[right] at (D) {$D$};
\node[above right,font=\small] at (M) {$M$};
\node[right] at (2,3) {$\mathcal{R}$};
\draw[dotted,-{Latex}] (0,-0.3) -- (4.3,-0.3) node[right]{$x$};
\draw[dotted,-{Latex}] (4,-0.6) -- (4,3.8) node[above]{$y$};
\node[below left,font=\small] at (4,-0.3) {$O$};
\end{tikzpicture}
\end{center}
\end{enumerate}

\vspace{6pt}
\nota{En las preguntas 4 a 5 complete los espacios en blanco o elija la (las) opción (opciones) correcta (correctas) según el caso. Tenga presente que los datos en cada literal son independientes y pueden variar. NOTA: En esta sección se califica sólo la respuesta y no se tiene en cuenta el procedimiento.}

\begin{enumerate}
\setcounter{enumi}{3}
\item{}[20pt] En la figura se observa el triángulo de vértices $O=\vect{0}{0}$, $Z$ y $W$. También se tiene la recta $\mathcal{L}$ que contiene a los puntos $Z$ y $W$ y otras siete rectas paralelas a $\mathcal{L}$, equidistantes entre sí, por ejemplo $V=\dfrac{6}{8}Z$. Finalmente, el vector $Z$ es ortogonal a la recta $\mathcal{L}$.

\begin{center}
\begin{tikzpicture}[scale=1]
\coordinate (O) at (4,0);
\coordinate (W) at (0,0.4);
\coordinate (Z) at (0.7,3.2);
\draw[thick] (O) -- (W) -- (Z) -- cycle;
\draw[dotted,-{Latex}] ($(W)!-0.15!(Z)$) -- ($(Z)!1.3!(W)$) node[above left]{$\mathcal{L}$};
\foreach \t in {0.15,0.3,0.45,0.6,0.75,0.9}{
  \draw[dashed] ($(O)!\t!(W)$) -- ($(O)!\t!(Z)$);
}
\node[font=\small] at ($(O)!0.15!(W)$) [below] {$U$};
\node[font=\small] at ($(O)!0.4!(W)$) [below] {$A$};
\node[font=\small] at ($(O)!0.7!(W)$) [below] {$C$};
\node[font=\small] at ($(O)!0.75!(Z)$) [right] {$B$};
\node[font=\small] at ($(O)!0.9!(Z)$) [above right] {$V$};
\node[left] at (W) {$W$};
\node[above] at (Z) {$Z$};
\node[right] at (O) {$O$};
\draw[dotted,-{Latex}] (0,-0.4) -- (0,3.6) node[above]{$y$};
\draw[dotted,-{Latex}] (-0.5,0) -- (4.6,0) node[right]{$x$};
\end{tikzpicture}
\end{center}

\begin{enumerate}[label=(\alph*)]
\item{}[5pt] $(B-C)\cdot(4Z)=$ \underline{\hspace{4cm}}.
\item{}[5pt] Si $W-Z=\vect{2\cos(4\pi/3)}{2\sin(4\pi/3)}$, entonces el ángulo $\alpha$ de inclinación de la recta $\mathcal{L}$ es $\alpha=$ \underline{\hspace{4cm}}.
\item{}[5pt] Si $A=tU+(1-t)C$, entonces $t=$ \underline{\hspace{4cm}}.
\item{}[5pt] Si $\|V\|=12$, calcular $V\cdot B=$ \underline{\hspace{4cm}}.
\end{enumerate}

\item{}[20pt] En la figura se observa un cuadrado cuyos vértices son los puntos $P$, $Q$, $R$ y $S$ y cuyos lados están sobre de las rectas $\mathcal{L}_1$, $\mathcal{L}_2$, $\mathcal{L}_3$ y $\mathcal{L}_4$. También se tiene una circunferencia $\mathcal{C}$ cuyo centro se encuentra en el origen $O=\vect{0}{0}$, que está inscrita en el cuadrado (es decir, es tangente a sus cuatro lados). Una de las diagonales del cuadrado pertenece a la recta $\mathcal{L}$.

\begin{center}
\begin{tikzpicture}[scale=1.1]
\draw (0,0) circle (1.3);
\coordinate (P) at (1.3,0.9);
\coordinate (Q) at (-0.9,1.3);
\coordinate (R) at (-1.3,-0.9);
\coordinate (S) at (0.9,-1.3);
\draw[thick] ($(Q)!-0.55!(P)$) -- ($(P)!-0.55!(Q)$);
\draw[thick] ($(R)!-0.55!(Q)$) -- ($(Q)!-0.55!(R)$);
\draw[thick] ($(S)!-0.55!(R)$) -- ($(R)!-0.55!(S)$);
\draw[thick] ($(P)!-0.55!(S)$) -- ($(S)!-0.55!(P)$);
\draw[thick,dashed] ($(R)!-0.5!(P)$) -- ($(P)!-0.5!(R)$);
\draw[dotted,-{Latex}] (0,-2.1) -- (0,2.1) node[above]{$y$};
\draw[dotted,-{Latex}] (-2.1,0) -- (2.3,0) node[right]{$x$};
\node[right,font=\small] at (0.1,-0.15) {$O$};
\node[above right] at (P) {$P$};
\node[above left] at (Q) {$Q$};
\node[below left] at (R) {$R$};
\node[below right] at (S) {$S$};
\node[left,font=\small] at (-0.7,0.7) {$\mathcal{C}$};
\node[font=\small] at (-2.0,-1.75) {$\mathcal{L}$};
\node[above,font=\small] at (0.2,2.05) {$\mathcal{L}_1$};
\node[left,font=\small] at (-2.05,0.2) {$\mathcal{L}_2$};
\node[below,font=\small] at (-0.2,-2.05) {$\mathcal{L}_3$};
\node[right,font=\small] at (2.05,-0.2) {$\mathcal{L}_4$};
\end{tikzpicture}
\end{center}

\begin{enumerate}[label=(\alph*)]
\item{}[5pt] Si la circunferencia tiene la ecuación $\mathcal{C}=\left\{X\in\mathbb{R}^2: \|X\|=\sqrt{\dfrac{10}{2}}\right\}$ entonces, la distancia entre $\mathcal{L}_2$ y $\mathcal{L}_4$ viene dada por $\operatorname{dist}(\mathcal{L}_2,\mathcal{L}_4)=$ \underline{\hspace{3cm}}.
\item{}[5pt] Si la pendiente de $\mathcal{L}_3$ es $m_3=-\dfrac{1}{2}$ entonces, la pendiente de $\mathcal{L}_2$ es $m_2=$ \underline{\hspace{3cm}}.
\item{}[5pt] Si $Q=\vect{-1}{3}$, un director de la recta $\mathcal{L}$ es $D=$ \underline{\hspace{3cm}}.
\item{}[5pt] Si $S=\vect{1}{-3}$ entonces el lado del cuadrado tiene longitud $\ell=$ \underline{\hspace{3cm}}.
\end{enumerate}

\end{enumerate}

\examfooter
\end{document}
